\centerline{\bf \Huge Начало Чётности II}

\begin{flushright}
        \scriptsize{Чётность не меняют только чётные числа - этим они и хороши. Начальник}
\end{flushright}

\problem{\textbf{1}} Во время турнира по квиддичу на Гарри Поттера напали 2025 дементоров. Когда он проскочил на метле мимо двух из них, они бросились на Гарри, промахнулись и превратились в камни. Гарри повторял этот маневр еще раз, и еще раз, и еще до тех пор, пока все дементоры не превратились в камни. Докажите, что Гарри Поттер что-то наколдовал.

\problem{\textbf{2}} Патрик Бейтман умножил сумму двух чисел на их произведение. Мог ли он получить в результате число 2025?

\problem{\textbf{3}} На доску размером 2024$\times$2024 пролили краску. Может ли количество испачканных клеток быть на 2025 меньше количества чистых клеток? 

\problem{\textbf{4}} Хулиган Герман порвал школьную стенгазету на 3 части. После этого он взял один из кусков и тоже порвал на 3 части. Потом опять один из кусков порвал на 3 части и т.д. Могло ли у него в итоге получиться 100 частей?

\problem{\textbf{5}} Можно ли все натуральные числа от 1 до 800 разбить на пары так,
чтобы сумма любой пары чисел делилась на 6? 

\problem{\textbf{6}} Как-то раз Дядя Фёдор, Матроскин и Шарик отправились
с почты домой. Дядя Фёдор вышел первым, а Матроскин последним. По дороге домой Дядя
Фёдор обгонял других, либо его обгоняли ровно 8 раз. Матроскин обгонял других, либо его
обгоняли ровно 6 раз. Известно, что Дядя Фёдор пришел домой позже, чем Шарик. В каком
порядке друзья пришли домой?

\problem{\textbf{7}} На чудо-дереве росли 2026 апельсинов и 2025 бананов. Каждый день садовник снимает с дерева ровно два фрукта. Если снятые фрукты одинаковы, то на дереве появляется новый
банан, а если разные — новый апельсин. В конце концов на дереве оказался только один фрукт.
Какой именно и почему?

